%%%%%%%%%%%%%%%%%%%% main.tex %%%%%%%%%%%%%%%%%%%%%%%%%%%%%%%%%%%
%
% Local Resistance-Based Design-Space Analysis of Polyamide Microtubes
% for Compact Heat Exchangers.
% Springer svmult contributed-volume manuscript (LNME, 155 mm x 235 mm).
%
% Build: pdflatex -> bibtex -> pdflatex x2  (or: latexmk -pdf main.tex)
%
% Section files (sections/):
%   00  Nomenclature
%   01  Introduction
%   02  Methods
%   03  Results
%   04  Case Study
%   05  Discussion
%   06  Conclusions
%   07  Acknowledgement
%   08  Appendix figures
%
%%%%%%%%%%%%%%%%%%%%%%%%%%%%%%%%%%%%%%%%%%%%%%%%%%%%%%%%%%%%%%%%%%

\documentclass[graybox]{svmult}
\pdfminorversion=7        % match MATLAB/Ghostscript vector plot exports
\pdfobjcompresslevel=0    % keep object streams uncompressed for parser compatibility

% Springer book format: 155 mm × 235 mm (LNME series).
% svmult.cls sets textwidth=117mm and textheight=191mm for this trim size
% but leaves the paper size unset. geometry is used here to set the PDF
% output page size and position the text block correctly; it overrides
% svmult's margin settings but preserves the text-area dimensions.
% Layout: inner=outer=19 mm, top-to-header=13 mm, bottom≈23 mm.
\usepackage[
  paperwidth  = 155mm,
  paperheight = 235mm,
  textwidth   = 117mm,
  textheight  = 191mm,
  headheight  = 12pt,
  headsep     = 12pt,
  top         = 13mm,
  inner       = 19mm,
  includehead
]{geometry}

% PDF text encoding / copy-paste support for accented names in pdflatex.
\usepackage{cmap}
\usepackage[T1]{fontenc}
\usepackage[utf8]{inputenc}

\IfFileExists{glyphtounicode.tex}{%
  \input{glyphtounicode}%
  \pdfgentounicode=1%
}{}

% Springer-style Times/Helvetica/Courier font setup.
\usepackage{mathptmx}
\usepackage{helvet}
\usepackage{courier}
\usepackage{type1cm}

\usepackage{makeidx}
\usepackage{graphicx}
\usepackage{multicol}
\usepackage[bottom]{footmisc}
\usepackage{textcmds}
\usepackage{amssymb,amsbsy,amsmath}
\usepackage[square,numbers,sort&compress]{natbib}
\usepackage{url}
\usepackage{subcaption}
\usepackage{booktabs}
\usepackage{array}
\usepackage{longtable}
\usepackage{rotating}
\usepackage{pdflscape}
\usepackage{needspace}
\usepackage{placeins}
\usepackage{tabularx}

\captionsetup{labelsep=space}

\newcommand{\DoeblerPrint}{D\"obler}
\newcommand{\be}{\begin{equation}}
\newcommand{\ee}{\end{equation}}
\newcommand{\ds}{\displaystyle}
\newcommand{\ETD}{\ensuremath{\mathrm{ETD}}}
\def\Cdot{\pmb{\cdot}}

% All figures are local to figures/ (self-contained submission package).
\graphicspath{{figures/}}

\makeindex

%%%%%%%%%%%%%%%%%%%%%%%%%%%%%%%%%%%%%%%%%%%%%%%%%%%%%%%%%%%%%%%%%%

\begin{document}

% Starred title follows the contributed-volume author style without chapter prefix.
\title*{Local Resistance-Based Design-Space Analysis\\
of Polyamide Microtubes\\
for Compact Heat Exchangers}
\titlerunning{Polyamide Microtube HX: Local Resistance Design-Space Analysis}

\author{
Philip \DoeblerPrint\inst{1} \and
Michael Henzler\inst{1} \and
Michael Auerbach\inst{1} \and
Andr\'e Casal Kulzer\inst{2}
}
\authorrunning{P. \DoeblerPrint{} et al.}

\institute{%
  \inst{1}Esslingen University of Applied Sciences,
  Kanalstra{\ss}e 33,
  73728 Esslingen am Neckar, Germany\\
  \email{philip.doebler@hs-esslingen.de}
  \at
  \inst{2}Institute of Automotive Engineering Stuttgart (IFS),
  University of Stuttgart,
  Pfaffenwaldring 12,
  70569 Stuttgart, Germany
}

\maketitle

\abstract{%
In automotive or aerospace contexts, air--liquid hollow-fiber heat exchangers (HFHX)
offer a high-area tube-bundle architecture for compact mobile cooling systems,
but metallic microtube bundles remain costly and manufacturing-intensive.
Polymer fibers may reduce tube-supply cost and mass, yet their low wall
conductivity, pressure integrity, and manufacturing tolerances must be evaluated
together rather than through surface-area arguments alone. In the sub-millimeter
range, the standard conductivity objection to polyamide (PA) loses quantitative
force: when outer convective resistance dominates, the PA wall contributes only
about $2$--$6\,\%$ of total thermal resistance at the material-specific minimum
wall ($t_{\mathrm{min},\mathrm{PA}}=0.025\,\mathrm{mm}$).
This paper tests whether that local compensation survives package and
manufacturing constraints.
A resistance-based design-space sweep over $(d_{\mathrm{o}},t)$ uses VDI Heat
Atlas G1/G7 correlations anchored to an aluminum benchmark and converts the
local heat-transfer coefficient to package-scaled conductance
$k_{\mathrm{o}}A_{\mathrm{o}}$.
Six constraints are applied to both materials: minimum wall thickness, coolant
throughput, straight-tube friction pressure drop, derated burst pressure,
capillary rise in the potting zone, and a preliminary tube-supply-cost cutoff.
The resulting design-boundary map shows that, under the normal extrusion
tolerance, a conditional sub-millimeter PA prototype target remains near
$d_{\mathrm{o}}\approx 0.35$--$0.65\,\mathrm{mm}$ at wall thicknesses
moderately above the process minimum.
The corresponding aluminum minimum-wall option is feasible only from
$d_{\mathrm{o}}\approx 0.63$--$0.65\,\mathrm{mm}$ because the example
tube-supply-cost cutoff excludes smaller aluminum diameters under the assumed
benchmark pitch.
The PA interval is therefore a model-derived prototype target that still
requires tests of wall tolerances, potting geometry, and bundle durability.
\keywords{hollow-fiber heat exchangers, polyamide microtubes,
tube-bank heat transfer, VDI Heat Atlas, design-space constraints,
potting capillarity}%
}

\clearpage
\input{sections/00_nomenclature}
\clearpage
\input{sections/01_introduction}
\clearpage
\input{sections/02_methods}
\section{Results}
\label{sec:results}

All results are derived from the analytical sweep of Section~\ref{sec:methods}
at fixed operating conditions ($u_{\mathrm{air},\infty} = 5\,\mathrm{m\,s^{-1}}$,
$v_{\mathrm{i}} = 0.5\,\mathrm{m\,s^{-1}}$). Their combination under one set of
constraint values is treated separately as a case study in
Section~\ref{sec:case-study}.

\subsection{Bundle-scale conductance map}
\label{sec:results-ko}

Figure~\ref{fig:ko} shows the package-scaled conductance
$k_{\mathrm{o}}A_{\mathrm{o}}$ over the $(d_{\mathrm{o}},\,\tau)$ plane for
aluminum and PA, with the respective minimum-wall constraints.
Both maps share a common color scale, so the absolute conductance gap is
directly visible.
Smaller outer diameters both raise $k_{\mathrm{o}}$ and increase the available
outer surface area at fixed package size; the resistance distribution behind
this trend is quantified in the next subsection.

\begin{figure}[htbp]
\centering
\includegraphics[width=\textwidth]{20_bundle_kA_alu_pa_portrait}
\caption{Package-scaled heat-transfer conductance
  $k_{\mathrm{o}}A_{\mathrm{o}}$ for aluminum and PA over the
  $(d_{\mathrm{o}},\tau)$ design space. Line and symbol conventions as defined
  in Sect.~\ref{sec:plot-conventions}; the $d_{\mathrm{o}}$ axis is logarithmic.}
\label{fig:ko}
\end{figure}

\pagebreak

\subsection{Thermal resistance distribution}
\label{sec:results-shares}

Figure~\ref{fig:shares} maps the three fractional resistances $\phi_{\mathrm{i}}$, $\phi_{\mathrm{w}}$, and $\phi_{\mathrm{o}}$
for aluminum and PA tubes over the $(d_{\mathrm{o}},\,\tau)$ design space.

\begin{figure}[htbp]
\centering
\includegraphics[width=\textwidth]{12_shares_alu_pa_portrait.pdf}
\caption{Fractional thermal resistance shares $\phi_{\mathrm{i}}$,
  $\phi_{\mathrm{w}}$, and $\phi_{\mathrm{o}}$ for Al and PA tubes. Line and
  symbol conventions as defined in Sect.~\ref{sec:plot-conventions}.}
\label{fig:shares}
\end{figure}

\pagebreak
The fractional resistance shares are defined as
\[
  \phi_j = \frac{R_j}{R_{\mathrm{tot}}}, \qquad j \in \{\mathrm{i},\,\mathrm{w},\,\mathrm{o}\},
\]
with the total resistance $R_{\mathrm{tot}} = R_{\mathrm{i}} + R_{\mathrm{w}} + R_{\mathrm{o}}$ and
$\phi_{\mathrm{i}} + \phi_{\mathrm{w}} + \phi_{\mathrm{o}} = 1$ by construction.

The outer convective resistance $\phi_{\mathrm{o}}$ constitutes the dominant fraction across
the small-diameter minimum-wall region.
Even for PA, where $\lambda_{\mathrm{w}} = 0.25\,\mathrm{W\,m^{-1}\,K^{-1}}$ represents an
880-fold reduction relative to aluminum, the wall resistance share $\phi_{\mathrm{w}}$
is $2.1$--$5.5\,\%$ along the PA minimum-wall line for
$d_{\mathrm{o}}=0.3$--$2.0\,\mathrm{mm}$.
When the outer convective resistance dominates, reducing wall conductivity by two orders
of magnitude attenuates $k_{\mathrm{o}}$ in proportion to $\phi_{\mathrm{w}}$ -- the penalty shrinks
with the wall-resistance share.
The corresponding outer share is $\phi_{\mathrm{o}}=83.8$--$85.4\,\%$ in this range.
At larger wall-thickness ratios, $\phi_{\mathrm{w}}$ grows and the geometry becomes less
attractive for compact cooling applications.

The inner convective share $\phi_{\mathrm{i}}$ is secondary but still relevant under the
conservative EGL5050 VDI G1 assumption; along the same line it
is $\phi_{\mathrm{i}}=9.6$--$13.2\,\%$, while the outer-side resistance remains
the largest term.



\FloatBarrier
\subsection{Coolant flow limits}
\label{sec:coolant-flow-limits}

Two coolant-side hydraulic limits are applied to the fixed-$v_{\mathrm{i}}$
sweep.
In Fig.~\ref{fig:coolant_throughput}, the geometrically similar pitch
assumption in Eq.~\eqref{eq:coolant_throughput} is shown. 

\begin{figure}[htbp]
\centering
\includegraphics[width=80mm]{19_coolant_throughput_lmin}
\caption{Coolant throughput diagnostic at fixed tube-side velocity. Line and
  symbol conventions as defined in Sect.~\ref{sec:plot-conventions}; the
  $d_{\mathrm{o}}$ axis is logarithmic.}
\label{fig:coolant_throughput}
\end{figure}

It makes the total flow depend
on the open-area ratio $(d_{\mathrm{i}}/d_{\mathrm{o}})^2$ under this
scaling, rather than on $d_{\mathrm{o}}$ alone.
At the fixed tube-side velocity used here, the increasing tube count at smaller
diameter is offset by the smaller flow area per tube; the associated wall-shear
effect is evaluated in the separate pressure-loss check below.

\Needspace{15\baselineskip}
\FloatBarrier
\subsection{Coolant pressure-loss limits}
\label{sec:coolant-pressure-loss-limits}

The pressure-loss map in Fig.~\ref{fig:dp_i_fric} follows the laminar
Hagen--Poiseuille scaling in Eq.~\eqref{eq:dp_i_laminar} over the PA
minimum-wall range.
Along the PA minimum-wall line, the affected region is confined to the smallest diameters:
for $d_{\mathrm{o}}\gtrsim 0.31\,\mathrm{mm}$ the straight-tube friction
pressure stays below $0.5\,\mathrm{bar}$.

\begin{figure}[htbp]
\centering
\includegraphics[width=80mm]{18_tube_side_friction_pressure_drop}
\caption{Straight-tube coolant pressure loss at fixed tube-side velocity. Line
  and symbol conventions as defined in Sect.~\ref{sec:plot-conventions}; the
  $d_{\mathrm{o}}$ axis is logarithmic.}
\label{fig:dp_i_fric}
\end{figure}

\FloatBarrier
\subsection{Tolerance-adjusted pressure-integrity check}
\label{sec:results-burst}

Figure~\ref{fig:burst} shows the tolerance-adjusted Lam\'e burst-pressure maps
for aluminum and PA over the $(d_{\mathrm{o}},\,\tau)$ plane.
Because $\sigma_{\mathrm{lim},\mathrm{Al}} / \sigma_{\mathrm{lim},\mathrm{PA}} = 20$, the PA isolines
are shifted strongly toward thicker walls compared with aluminum at the same
internal pressure.
In the thermally attractive small-diameter range, however, parts of the PA
window remain above typical automotive cooling-circuit operating pressures.
This is only a diagnostic observation; the integrated design-boundary map uses
the stricter criterion
$p_{\mathrm{b},\mathrm{tol}}\ge 6\,\mathrm{bar}$.
The comparison between rows shows the sensitivity to absolute tolerance:
at the PA minimum wall, the normal route leaves
$t_{\mathrm{loc},\mathrm{min}}=5\,\mu\mathrm{m}$, whereas the high-precision route leaves
$20\,\mu\mathrm{m}$.
At fixed $t_{\mathrm{min}}$ and $\Delta t$, smaller $d_{\mathrm{o}}$ increases
$(t_{\mathrm{min}}-\Delta t)/d_{\mathrm{o}}$ and raises the Lam\'e geometry factor.
The check shifts attention from isolated tube-wall burst to
bundle-level mechanisms outside this idealized calculation.

\begin{figure}[htbp]
\centering
\includegraphics[width=\textwidth]{15_burst_tolerance_alu_pa_portrait}
\caption{Tolerance-adjusted Lam\'e burst pressure
  $p_{\mathrm{b},\mathrm{tol}}$ for aluminum and PA tubes.
  Columns compare materials; rows compare
  $\Delta t=0.020\,\mathrm{mm}$ and $\Delta t=0.005\,\mathrm{mm}$. Line and
  symbol conventions as defined in Sect.~\ref{sec:plot-conventions}; the
  $d_{\mathrm{o}}$ axis is logarithmic.}
\label{fig:burst}
\end{figure}

\FloatBarrier
\subsection{Capillary-rise process constraint}
\label{sec:results-capillary-rise}

Figure~\ref{fig:capillary_rise} maps the predicted resin rise for aluminum and
PA columns under effective acceleration factors $G=1$, 5, and 10.
At $G=1$, sub-millimeter clear gaps generate process-relevant resin rise;
because Eq.~\eqref{eq:capillary_rise} scales as $1/G$, increasing
acceleration to $G=5$--$10$ reduces the predicted rise by the same factor.
For the benchmark pitch ratio $S_{\mathrm{L}}/\!d_{\mathrm{o}}=2.00$, the
$h_{\mathrm{cap}}=2\,\mathrm{mm}$ line at $G=10$ is near
$d_{\mathrm{o}}\approx 0.31\,\mathrm{mm}$ for the aluminum estimate and
$d_{\mathrm{o}}\approx 0.27\,\mathrm{mm}$ for the PA\,/\,process proxy. The two
capillary limits are much closer than a PA material-constant
interpretation would suggest. Compact pitch geometries that are thermally
attractive still require a controlled potting route, such as centrifugal
potting, as an independent process constraint.

\begin{figure}[htbp]
\centering
\includegraphics[width=\textwidth]{14_capillary_rise_alu_pa_portrait}
\caption{Predicted capillary rise $h_{\mathrm{cap}}$ for aluminum and PA
  potting gaps under effective acceleration factors $G=1$, 5, and 10. Line and
  symbol conventions as defined in Sect.~\ref{sec:plot-conventions}; the
  $d_{\mathrm{o}}$ axis is logarithmic.}
\label{fig:capillary_rise}
\end{figure}

\FloatBarrier

\clearpage
\section{Case study: baseline prototype-target scenario}
\label{sec:case-study}

The preceding model results describe the individual thermal, hydraulic,
pressure-integrity, and potting-process relationships within their stated
limits. This case study combines them with the tube-supply-cost criterion using
one internally consistent set of example values. The values in
Table~\ref{tab:design_boundary_assumptions} define this example; they are not
design allowables or universal feasibility limits.

\subsection{Example design-boundary map}
\label{sec:case-design-boundary}

Figure~\ref{fig:design_boundary} combines the preceding thermal, hydraulic,
mechanical, process, and tube-supply-cost constraints in one example design-space
view. The baseline values and thresholds are summarized in
Table~\ref{tab:design_boundary_assumptions}.

\begin{table}[htbp]
\caption{Baseline thresholds used for the design-boundary map.}
\label{tab:design_boundary_assumptions}
\centering
\small
\setlength{\tabcolsep}{3pt}
\renewcommand{\arraystretch}{0.95}
\begin{tabular*}{\textwidth}{@{}ll@{\extracolsep{\fill}}ccc@{}}
\toprule
Boundary & Scope & Symbol & Limit & Unit \\
\midrule
Min. wall thickness & Al, drawn tube
  & $t$ & $\ge 0.070$ & $\mathrm{mm}$ \\
Min. wall thickness & PA, precision route
  & $t$ & $\ge 0.025$ & $\mathrm{mm}$ \\
Coolant throughput & fixed velocity model
  & $\dot{V}_{\mathrm{c}}$ & $\ge 10$ & $\mathrm{L\,min^{-1}}$ \\
Pressure drop & fixed friction model
  & $\Delta p_{\mathrm{i},\mathrm{fric}}$ & $\le 0.5$ & $\mathrm{bar}$ \\
Tube-supply cost & cost index
  & $C_{\mathrm{tube}}/C_{\mathrm{fin}}$ & $<5$ & -- \\
Burst pressure & normal extrusion
  & $p_{\mathrm{b},\mathrm{tol}}$ & $\ge 6$ & $\mathrm{bar}$ \\
Capillary rise & potting, $G=10$
  & $h_{\mathrm{cap}}$ & $\le 2$ & $\mathrm{mm}$ \\
\bottomrule
\end{tabular*}
\end{table}
\FloatBarrier

Under this example value set, cost, pressure-integrity, potting, and
coolant-side constraints remove large parts of the thermally attractive design
space. The remaining PA region defines a prototype target window for this value
set; turning it into a design rule would require data over production
tolerances, potting quality, and pressure cycling.
Under the example normal-route tolerance, the PA minimum-wall line itself
fails the $6\,\mathrm{bar}$ burst criterion:
$t_{\mathrm{min},\mathrm{PA}}=0.025\,\mathrm{mm}$ leaves a minimum local wall of
$t_{\mathrm{loc},\mathrm{min}}=5\,\mu\mathrm{m}$.
The PA region accepted by the combined constraints requires walls moderately
above the example process minimum,
$t \gtrsim \Delta t + 0.03\,d_{\mathrm{o}}$, about
$0.030$--$0.040\,\mathrm{mm}$ at sub-millimeter diameters.
Its lower diameter edge in this scenario is not set by the capillary limit.
Under the example $G=10$ benchmark-pitch convention and the
$h_{\mathrm{cap}}\le2\,\mathrm{mm}$ limit, the updated PA/process capillary
limit reaches its threshold near $d_{\mathrm{o}}\approx 0.27\,\mathrm{mm}$,
below the tube-supply-cost cutoff near $d_{\mathrm{o}}\approx 0.35\,\mathrm{mm}$
and the straight-tube friction limit near $d_{\mathrm{o}}\approx
0.31\,\mathrm{mm}$.
For aluminum, the case-study minimum-wall boundary starts at
$d_{\mathrm{o}}\approx 0.63$--$0.65\,\mathrm{mm}$ because the example
tube-supply-cost cutoff excludes smaller aluminum diameters.

\begin{figure}[htbp]
\centering
\includegraphics[width=\textwidth]{20_design_boundary_lines_alu_pa}
\caption{Design-boundary map for Al and PA tubes.
  Color encodes $k_{\mathrm{o}}A_{\mathrm{o}}$ for constraint-admitted points;
  boundary colors follow the legend and hatching marks excluded sides. Blue
  cross: validated aluminum benchmark. The $d_{\mathrm{o}}$ axis is logarithmic.}
\label{fig:design_boundary}
\end{figure}

\FloatBarrier
\subsection{Feasible PA-to-aluminum conductance comparison}
\label{sec:case-delta}

The comparison in Fig.~\ref{fig:tech_adj} uses the combined feasibility masks
from Table~\ref{tab:design_boundary_assumptions}. For
$m\in\{\mathrm{Al},\mathrm{PA}\}$, the mask
$\mathcal{F}_m(d_{\mathrm{o}},\tau)\in\{0,1\}$ is set to 1 when the design
point satisfies all applicable thresholds in
Table~\ref{tab:design_boundary_assumptions}, and to 0 otherwise.
For each feasible PA point, the aluminum comparator is the nearest feasible
aluminum geometry at the same wall-thickness ratio $\tau$ and at an equal or
larger outer diameter. 

\begin{figure}[htbp]
\centering
\includegraphics[width=\textwidth]{22_tech_adjusted_kA_change}
\caption{Case-study relative conductance difference
  $\Delta (k_{\mathrm{o}}A_{\mathrm{o}})^{\,\mathrm{feas}}$.
  Color marks PA points passing all example constraints.
  Each PA point is compared with the nearest constraint-admitted aluminum diameter at
  the same $\tau$; uncolored areas are excluded or have no admitted aluminum
  comparator. Solid black curves are isolines of
  $\Delta (k_{\mathrm{o}}A_{\mathrm{o}})^{\,\mathrm{feas}}$.
  Dashed dark-blue and dotted dark-green curves mark the Al and PA combined-constraint
  lower-diameter boundaries, respectively. Blue cross: validated aluminum
  benchmark geometry. The $d_{\mathrm{o}}$ axis is logarithmic.}
\label{fig:tech_adj}
\end{figure}

The feasible PA-to-aluminum comparison is defined by
Eq.~\eqref{eq:feasible_pa_al_comparison},
\be
  \begin{aligned}
  d_{\mathrm{o},\mathrm{Al}}^{*}(d_{\mathrm{o}},\tau)
  &=
  \min\left\{
    \tilde d_{\mathrm{o}}\ge d_{\mathrm{o}}\; \middle|\;
    \mathcal{F}_{\mathrm{Al}}(\tilde d_{\mathrm{o}},\tau)=1
  \right\},\\
  \Delta (k_{\mathrm{o}}A_{\mathrm{o}})^{\,\mathrm{feas}}
  &=
  \frac{
    (k_{\mathrm{o}}A_{\mathrm{o}})_{\mathrm{PA}}(d_{\mathrm{o}},\tau)
  }{
    (k_{\mathrm{o}}A_{\mathrm{o}})_{\mathrm{Al}}
    (d_{\mathrm{o},\mathrm{Al}}^{*},\tau)
  }
  - 1,
  \qquad
  \mathcal{F}_{\mathrm{PA}}(d_{\mathrm{o}},\tau)=1 .
  \end{aligned}
  \label{eq:feasible_pa_al_comparison}
\ee

In this map a positive value marks a PA point that exceeds the nearest
constraint-admitted aluminum option at the same wall-thickness ratio, whereas a
negative value marks a point where aluminum remains superior. The relevant PA
target is the conditional sub-millimeter interval identified in
Sect.~\ref{sec:case-design-boundary}, where no constraint-admitted aluminum
comparator exists at a smaller diameter. The comparison isolates
geometry-dependent conductance; burst/leak tests and bundle heat-rate and
pressure-drop tests would have to show prototype qualification and system
performance.

\FloatBarrier

\section{Discussion}
\label{sec:discussion}

\subsection{Meaning and limits of the PA target interval}

The results define a conditional PA prototype target interval without
establishing a general PA thermal advantage. In the small-diameter, thin-wall
regime, the outer convective resistance dominates, so the low PA wall conductivity affects
$k_{\mathrm{o}}$ mainly through the wall-resistance share
$\phi_{\mathrm{w}}$. This is consistent with the near-matched
PA612--stainless-steel comparison~\cite{Astrouski2026}, where a staggered
$d_{\mathrm{o}}=1.23\,\mathrm{mm}$, $t=0.12\,\mathrm{mm}$ PA612 bank at
$2$--$10\,\mathrm{m\,s^{-1}}$ lost ${\approx}\,7\,\%$ of mean heat rate; its
staggered layout, thicker relative wall, and steel comparator test the
resistance-share mechanism for a different geometry, whereas the present map
uses inline benchmark pitch and the combined constraints.
Positive values in Fig.~\ref{fig:tech_adj} mean that a feasible PA
point exceeds the nearest feasible aluminum alternative under the constraints
used in Fig.~\ref{fig:design_boundary}.
The PA interval depends on more than pure area scaling; it emerges from the
combined effect of outer-area conductance, resistance shares, the fixed
packing assumptions, and the applied tolerance, potting, burst, and cost
limits.

This interval is bounded by manufacturing and durability constraints. The PA
region depends on wall thicknesses above the nominal process minimum under the
normal tolerance case, controlled potting, and the derated burst-strength
input.
Although the lower PA density favors tube mass, assembly-level mass and
performance can still be dominated by potting, supports, tube sheets, bypass
flow, and pressure integrity.

\Needspace{15\baselineskip}
\subsection{Manufacturing and durability constraints}

Wall tolerance is the first practical limit. At
$d_{\mathrm{o}}=0.5\,\mathrm{mm}$ and $t=0.025\,\mathrm{mm}$, an eccentricity of
$\pm 5\,\mu\mathrm{m}$ is already a $\pm 20\,\%$ wall-thickness variation.
For the normal route used in the integrated map,
$\Delta t=20\,\mu\mathrm{m}$ leaves a
$5\,\mu\mathrm{m}$ local wall at the PA process minimum; the burst
criterion forces walls above $t_{\mathrm{min},\mathrm{PA}}$. A Lam\'e pressure
check at nominal wall thickness would be non-conservative. Thin-walled
sub-millimeter PA hollow fibers have been demonstrated in heat-exchanger
applications~\cite{Kroulikova2021,Bulejko2022}, but tolerances, concentricity,
continuous length, tube-count yield, and cost at a given diameter--wall
combination remain process- and supplier-specific~\cite{Mount2017}; public
supplier capability data do not document the full combination assumed here.

Potting is the second practical limit. The capillary-rise limit in
Fig.~\ref{fig:design_boundary} is a potting-process constraint: resin rise
shortens active length but does not change the local heat-transfer
correlations. Centrifugal potting addresses the body-force term in the
capillary balance, so membrane-module practice gives useful process guidance.
For this heat-exchanger geometry, however, the $G=10$ assumption still has to be
checked with the final tube surface, resin system, gel time, fixture, and
potting depth.

Pressure integrity must be assessed at the supported potted bundle level;
isolated tube-wall burst covers a single failure mode. In early prototype
builds, failures occurred in the frame, potting, and tube--potting interface,
with bundle deformation contributing to the failure path. The $6\,\mathrm{bar}$ Lam\'e
criterion excludes tube-wall geometries with inadequate idealized pressure
margin. Supports, creep, aging, fatigue, pressure cycling, moisture uptake and
hydrolysis in the glycol--water coolant, and the adhesion and long-term
durability of the tube--potting interface are bundle qualification items.
Published pulsating-pressure fatigue
tests~\cite{Brozova2016} and temperature-dependent burst
measurements~\cite{Bulejko2022} of polymeric hollow-fiber heat-transfer
surfaces provide first experimental anchors for these mechanisms; the present
bundle still needs its own support, potting, and coolant-exposure tests.
Higher-temperature
polymers such as polyether ether
ketone (PEEK) may improve these margins but require their own material/process
sweep and tests~\cite{Peters2017}.

\subsection{Model scope}

The heat-transfer model is a VDI-correlation sweep at one operating point with
constant fluid properties and benchmark inline pitch
ratios~\cite{GnielinskiG7,GnielinskiG1}; the resulting
$k_{\mathrm{o}}A_{\mathrm{o}}$ gives a package-scaled conductance potential.
Full heat-rate and vehicle-system predictions require heat-rate and
outer- and tube-side pressure-drop evaluation including
headers and bypass, fan and pump power, and tests on the realized bundle. The
capillary map informs potting trials, and the tube-supply-cost index flags
diameter ranges for supplier quotation rather than setting a production cost.

\section{Conclusions}
\label{sec:conclusions}

An analytical design-space sweep maps conductance, resistance shares, and six
feasibility criteria for aluminum and PA tubes.
The sweep is anchored to the measured aluminum benchmark from the companion
study~\cite{Henzler2026}; PA points are extrapolated from that basis.
All maps hold at the fixed operating point
($u_{\mathrm{air},\infty}=5\,\mathrm{m\,s^{-1}}$,
$v_{\mathrm{i}}=0.5\,\mathrm{m\,s^{-1}}$) with the VDI G1/G7 correlations; they
are local correlation results for this operating window.
The conclusions are:

\begin{enumerate}

\item In the small-diameter PA minimum-wall interval
  ($d_{\mathrm{o}}=0.3$--$2.0\,\mathrm{mm}$), Fig.~\ref{fig:shares} shows that
  the wall remains a minor resistance term and outer convection remains the
  controlling resistance term.
  The conductivity penalty is governed by the fractional wall-resistance share,
  so the 880-fold conductivity difference to aluminum enters through a small
  resistance fraction.

\item With the six criteria applied under the normal-route tolerance
  ($\Delta t=0.020\,\mathrm{mm}$), the design-boundary map places the
  case-study PA target near $d_{\mathrm{o}}\approx 0.35\,\mathrm{mm}$, where
  the exploratory tube-supply-cost cutoff is more restrictive than the updated
  PA/process capillary limit under the $G=10$,
  $h_{\mathrm{cap}}\le2\,\mathrm{mm}$ case (Fig.~\ref{fig:design_boundary}).
  At sub-millimeter diameters, this interval requires wall thicknesses
  moderately above the process minimum, about
  $0.030$--$0.040\,\mathrm{mm}$, to pass the $6\,\mathrm{bar}$ burst criterion.
  The aluminum minimum-wall option starts at
  $d_{\mathrm{o}}\approx 0.63$--$0.65\,\mathrm{mm}$, because the
  tube-supply-cost cutoff excludes smaller aluminum diameters under the
  benchmark-pitch assumptions.
  Admission of PA without an equal-diameter aluminum comparator remains
  conditional.

\item Because the identified PA interval depends on controlled potting at
  $G=10$, a derated PA burst-strength input, and wall thicknesses above the
  process minimum under the normal $\pm 20\,\mu\mathrm{m}$ extrusion tolerance,
  the practical design limit shifts from thermal resistance to manufacturing,
  pressure integrity, and durability.
  Early prototype builds shifted the failure concern from isolated tube burst
  toward frame stiffness, bundle deformation, and tube--potting integrity.
  The next qualification step is bundle-level structural testing of
  wall-tolerance capability, potting adhesion, pressure integrity, and
  long-term coolant-interface durability.

\end{enumerate}


\begin{acknowledgement}
The authors acknowledge the Esslingen Formula Student team for prototype
manufacturing and test-bench work.
Large Language Models were limited to AI-assisted copy editing of readability,
grammar, and formatting; autonomous content creation, data analysis, and
scientific conclusions were performed by the authors.
The authors reviewed all edits and take full responsibility for the final manuscript.

\end{acknowledgement}
\clearpage
\appendix
\section*{Appendix: Supplementary Figures and Photographs}
\label{sec:appendix_photographs}

\setcounter{figure}{0}
\renewcommand{\thefigure}{A.\arabic{figure}}

\begin{figure}[!htbp]
\centering
\includegraphics[width=0.5\textwidth]{capillary_rise.png}
\caption{Local resin rise between closely spaced vertical tubes.}
\label{fig:appendix_close_tube_capillary}
\end{figure}

\begin{figure}[!htbp]
\centering
\includegraphics[width=0.9\textwidth]{wedge_gap_resin_rise.pdf}
\caption{Wedge-gap resin-rise experiment used to estimate the effective
  full-gap capillary parameter
  \(C_{\mathrm{cap}} = 5.4 \pm 0.5\,\mathrm{mm^2}\) for the
  unfilled WEVOPOX/WEVODUR process configuration at \(22^\circ\mathrm{C}\).
  (\textbf{a}), (\textbf{b})~Front views of the two wedge plates. Dashed
  lines: scribed \(10\,\mathrm{mm}\) reference grid of each plate; numbers
  give the height above the scribed baseline. Blue dotted curves: joint fit
  of the digitized resin fronts. (\textbf{c})~Side view with the measured
  wedge angle \(\varepsilon\). The transparent Plexiglas/PMMA wedge provides
  optical access as a process proxy. Direct PA wetting measurements are
  outside this experiment.}
\label{fig:appendix_wedge_resin}
\end{figure}

\begin{figure}[htbp]
\centering
\includegraphics[width=0.8\textwidth]{manufactured_prototypes.jpg}
\caption{Manufactured polymeric prototypes.}
\label{fig:appendix_prototype_collection}
\end{figure}

\FloatBarrier

\bibliographystyle{spmpsci}
\bibliography{references}

\end{document}
